% R&R Other Project Information, Field 10: Facilities & Other Resources
% (NOFO N0001426SF004, p. 20). Adapted from the FLAIME DOE proposal,
% Appendix 2. Build: xelatex field10-facilities.tex
\documentclass[12pt,letterpaper]{article}
\usepackage[margin=1in]{geometry}
\usepackage{fontspec}
\IfFontExistsTF{Times New Roman}{\setmainfont{Times New Roman}}{\setmainfont{TeX Gyre Termes}}
\usepackage{microtype}
\usepackage{xcolor}
\usepackage{xurl}
\usepackage[colorlinks=true,urlcolor=blue!70!black]{hyperref}
\setlength{\parindent}{0pt}
\setlength{\parskip}{5pt}
\pagestyle{plain}
\newcommand{\todo}[1]{\textcolor{red}{[TODO Sterling: #1]}}

\begin{document}
\begin{center}\textbf{Facilities and Other Resources}\\
ONR FY2027 YIP, NOFO N0001426SF004. PI: Sterling G.\ Baird, Brigham Young University\end{center}

\textbf{PI laboratory, BYU Department of Mechanical Engineering.} A powder
metallurgy laboratory \todo{square footage and room} holds an AMAZEMET
rePOWDER ultrasonic atomizer with an induction crucible module, which makes
aluminum alloy powder from elemental and master-alloy feedstock weighed into
machined crucibles of high-purity aluminum, four compositions per week. The
room has a canopy hood, a fume hood, an industrial dehumidifier, dry powder
storage, and electrical, gas, compressed air, and chilled water services,
installed with BYU Environmental Health and Safety. The laboratory connects
directly to the BYU Microscopy Center by a door with no public access.

\textbf{BYU Microscopy Center.} Scanning electron microscopes with
energy-dispersive X-ray spectroscopy (EDS) and electron backscatter
diffraction (EBSD) (Apreo C, Helios 600 with a focused ion beam, and Verios), and a
Tecnai TF-20 transmission electron microscope. Instruments are shared and
have high availability.

\textbf{Other BYU facilities.} Metallographic preparation (PX500 automatic
and PX300 manual polishers, vibratory polishers, fume hood for etching, and
an electropolisher); a Keyence VHX-7000 digital optical microscope; the BYU
X-ray Diffraction Facility; the BYU Spectroscopic Facility (Agilent 8900
ICP-MS/MS for composition, and XPS); and the BYU Prototyping Lab (CNC
machining for crucibles and test specimens). Tensile tests use
\todo{name the BYU load frame or test service}.

\textbf{Computing.} The BYU Office of Research Computing operates a shared
cluster of about 28,000 CPU cores and 367 GPUs on 616 nodes
(\url{https://rc.byu.edu}) \todo{update the counts if they have changed
since April 2026}, available at no direct cost through fair-share
scheduling. CALPHAD calculations use Thermo-Calc \todo{confirm the BYU
license and the aluminum databases} and pycalphad.

\textbf{Laser powder bed fusion.} Builds are bought on a purchase order at
\$800 per build from a University of Utah user facility \todo{facility name,
printer model, and build envelope; the 2026 DOE proposal named an Aconity
MIDI with a 170~mm build plate and a 200 to 1000~W fiber laser}.

\textbf{Corrosion testing.} The electrochemical screen uses a potentiostat
requested in Year 1 (see Field 11). Neutral salt-spray exposures are bought
from an outside test laboratory \todo{name the laboratory once quoted}.

\textbf{Government-owned facilities.} None are required \todo{revise if a
Navy laboratory agrees to run stress corrosion tests on final candidates}.

\end{document}
